\documentclass[10pt,a4paper]{amsart}
\usepackage{amsmath,amssymb,amsthm,graphicx}
% Retained preamble for the Stacks Project chapter {CH} (shared project preamble).
% The chapter file pinned-original-source.tex only carries \input{{preamble}} on its first line,
% so no preamble text lives inside the chapter itself.  The Stacks Project uses one single shared
% preamble for every chapter: the snapshot stages it three times (sources/stacks-chow/preamble.tex,
% sources/stacks-etale-cohomology/preamble.tex and focused-seeds/AGP/sources/stacks-snapshot-20261011/
% preamble.tex) and those copies are identical in size.  The block below is copied mechanically,
% line for line, from sources/stacks-chow/preamble.tex lines 14-18 and 154-217: the Xy-pic setup and
% the native theorem declarations together with the mathematical macros.  Nothing was added, removed
% or redefined; the wrapper class and the amsmath/amssymb/amsthm/graphicx base packages are supplied
% by the assembler.
% For commutative diagrams we use Xy-pic
% Chapter: stacks-cohomology.tex -- Cohomology of Sheaves.
\usepackage[all]{xy}

% We use 2cell for 2-commutative diagrams.
\xyoption{2cell}

% Theorem environments.
%
\theoremstyle{plain}
\newtheorem{theorem}[subsection]{Theorem}
\newtheorem{proposition}[subsection]{Proposition}
\newtheorem{lemma}[subsection]{Lemma}

\theoremstyle{definition}
\newtheorem{definition}[subsection]{Definition}
\newtheorem{example}[subsection]{Example}
\newtheorem{exercise}[subsection]{Exercise}
\newtheorem{situation}[subsection]{Situation}

\theoremstyle{remark}
\newtheorem{remark}[subsection]{Remark}
\newtheorem{remarks}[subsection]{Remarks}

\numberwithin{equation}{subsection}

% Macros
%
\def\lim{\mathop{\mathrm{lim}}\nolimits}
\def\colim{\mathop{\mathrm{colim}}\nolimits}
\def\Spec{\mathop{\mathrm{Spec}}}
\def\Hom{\mathop{\mathrm{Hom}}\nolimits}
\def\Ext{\mathop{\mathrm{Ext}}\nolimits}
\def\SheafHom{\mathop{\mathcal{H}\!\mathit{om}}\nolimits}
\def\SheafExt{\mathop{\mathcal{E}\!\mathit{xt}}\nolimits}
\def\Sch{\mathit{Sch}}
\def\Mor{\mathop{\mathrm{Mor}}\nolimits}
\def\Ob{\mathop{\mathrm{Ob}}\nolimits}
\def\Sh{\mathop{\mathit{Sh}}\nolimits}
\def\NL{\mathop{N\!L}\nolimits}
\def\CH{\mathop{\mathrm{CH}}\nolimits}
\def\proetale{{pro\text{-}\acute{e}tale}}
\def\etale{{\acute{e}tale}}
\def\QCoh{\mathit{QCoh}}
\def\Ker{\mathop{\mathrm{Ker}}}
\def\Im{\mathop{\mathrm{Im}}}
\def\Coker{\mathop{\mathrm{Coker}}}
\def\Coim{\mathop{\mathrm{Coim}}}

% Boxtimes
%
\DeclareMathSymbol{\boxtimes}{\mathbin}{AMSa}{"02}

%
% Macros for moduli stacks/spaces
%
\def\QCohstack{\mathcal{QC}\!\mathit{oh}}
\def\Cohstack{\mathcal{C}\!\mathit{oh}}
\def\Spacesstack{\mathcal{S}\!\mathit{paces}}
\def\Quotfunctor{\mathrm{Quot}}
\def\Hilbfunctor{\mathrm{Hilb}}
\def\Curvesstack{\mathcal{C}\!\mathit{urves}}
\def\Polarizedstack{\mathcal{P}\!\mathit{olarized}}
\def\Complexesstack{\mathcal{C}\!\mathit{omplexes}}
% \Pic is the operator that assigns to X its picard group, usage \Pic(X)
% \Picardstack_{X/B} denotes the Picard stack of X over B
% \Picardfunctor_{X/B} denotes the Picard functor of X over B
\def\Pic{\mathop{\mathrm{Pic}}\nolimits}
\def\Picardstack{\mathcal{P}\!\mathit{ic}}
\def\Picardfunctor{\mathrm{Pic}}
\def\Deformationcategory{\mathcal{D}\!\mathit{ef}}

% --- appended 2026-10-12 (Chief): the project own definitions of the reference/slogan/history
% display environments, copied mechanically from the shared Stacks preamble lines 8-12
% (\usepackage{verbatim} plus the three \newenvironment lines, which the chapter preamble keeps
% verbatim from the project preamble). Nothing added, removed or redefined: these environments
% carry no mathematical content and the project itself typesets them as comments. Without these
% three lines a counted statement whose own text uses them fails with Environment slogan undefined.
\usepackage{verbatim}
\newenvironment{reference}{\comment}{\endcomment}
\newenvironment{slogan}{\comment}{\endcomment}
\newenvironment{history}{\comment}{\endcomment}
\usepackage{cleveref}
\crefname{originalreference}{original reference}{original references}
\Crefname{originalreference}{Original reference}{Original references}
\makeatletter
\newlabel{situation-locally-given}{{0D67}{source}{}{sourcepointer-1}{}}
\newlabel{situation-locally-given@cref}{{[originalreference][0][]0D67}{[1][0][]source}{}{}{}}
\makeatother
\crefname{originallemma}{original lemma}{original lemmas}
\Crefname{originallemma}{Original lemma}{Original lemmas}
\makeatletter
\newlabel{lemma-uniqueness}{{0D68}{source}{}{sourcepointer-2}{}}
\newlabel{lemma-uniqueness@cref}{{[originallemma][0][]0D68}{[1][0][]source}{}{}{}}
\makeatother
\crefname{originallemma}{original lemma}{original lemmas}
\Crefname{originallemma}{Original lemma}{Original lemmas}
\makeatletter
\newlabel{lemma-solution-in-finite-case}{{0D6A}{source}{}{sourcepointer-3}{}}
\newlabel{lemma-solution-in-finite-case@cref}{{[originallemma][0][]0D6A}{[1][0][]source}{}{}{}}
\makeatother
\crefname{originallemma}{original lemma}{original lemmas}
\Crefname{originallemma}{Original lemma}{Original lemmas}
\makeatletter
\newlabel{lemma-glueing-increasing-union}{{0D6B}{source}{}{sourcepointer-4}{}}
\newlabel{lemma-glueing-increasing-union@cref}{{[originallemma][0][]0D6B}{[1][0][]source}{}{}{}}
\makeatother
\title{BBD gluing lemma (general form, without boundedness)}
\author{The Stacks Project Authors}
\date{stacks-project:cohomology.tex:local-snapshot-20261011; GNU FDL 1.2 or later}
\begin{document}
\maketitle
\noindent\textit{Editorial note.} All delivered mathematical text is the author's own, copied verbatim from the pinned Stacks Project chapter file cohomology.tex; no statement or proof step was rewritten, reordered or removed.\par\medskip
\noindent\textit{Editorial note.} The counted span cites the lemmas recorded in assembly.external\_references with canonical tags as navigation pointers.\par\medskip
\noindent\textit{Allowed standard prerequisites (editorial).}\begin{itemize}
\item Uniqueness of solutions. Cohomology of Sheaves, Lemma cohomology-lemma-uniqueness. Solutions are unique up to unique isomorphism and have vanishing negative exts.
\item Finite glueing. Cohomology of Sheaves, Lemma cohomology-lemma-solution-in-finite-case. Existence for systems on a space with a finite covering.
\item Glueing along increasing unions. Cohomology of Sheaves, Lemma cohomology-lemma-glueing-increasing-union. Existence along well-ordered increasing unions with compatible isomorphisms.
\end{itemize}\medskip
\ifdefined\hypertarget\hypertarget{sourcepointer-1}{}\fi
\noindent\textit{Source reference.} Original reference 0D67: sources/stacks-chow/tags.txt, source lines 16512--16512. Cross-chapter standard prerequisite cited by the author inside the counted span; it is used as a pointer and its own proof is not part of this item.\par\medskip
\ifdefined\hypertarget\hypertarget{sourcepointer-2}{}\fi
\noindent\textit{Source reference.} Original lemma 0D68: sources/stacks-chow/tags.txt, source lines 16513--16513. Cross-chapter standard prerequisite cited by the author inside the counted span; it is used as a pointer and its own proof is not part of this item.\par\medskip
\ifdefined\hypertarget\hypertarget{sourcepointer-3}{}\fi
\noindent\textit{Source reference.} Original lemma 0D6A: sources/stacks-chow/tags.txt, source lines 16515--16515. Cross-chapter standard prerequisite cited by the author inside the counted span; it is used as a pointer and its own proof is not part of this item.\par\medskip
\ifdefined\hypertarget\hypertarget{sourcepointer-4}{}\fi
\noindent\textit{Source reference.} Original lemma 0D6B: sources/stacks-chow/tags.txt, source lines 16516--16516. Cross-chapter standard prerequisite cited by the author inside the counted span; it is used as a pointer and its own proof is not part of this item.\par\medskip
\begin{theorem}[BBD gluing lemma]
\label{theorem-glueing-bbd-general}
\begin{reference}
Special case of \cite[Theorem 3.2.4]{BBD}
without boundedness assumption.
\end{reference}
In Situation \ref{situation-locally-given} assume
\begin{enumerate}
\item $X = \bigcup_{U \in \mathcal{B}} U$,
\item for $U, V \in \mathcal{B}$ we have
$U \cap V = \bigcup_{W \in \mathcal{B}, W \subset U \cap V} W$,
\item for any $U \in \mathcal{B}$ we have $\Ext^i(K_U, K_U) = 0$
for $i < 0$.
\end{enumerate}
Then there exists an object $K$ of $D(\mathcal{O}_X)$
and isomorphisms $\rho_U : K|_U \to K_U$ in $D(\mathcal{O}_U)$ for
$U \in \mathcal{B}$ such that $\rho^U_V \circ \rho_U|_V = \rho_V$
for all $V \subset U$ with $U, V \in \mathcal{B}$.
The pair $(K, \rho_U)$ is unique up to unique isomorphism.
\end{theorem}

\begin{proof}
A pair $(K, \rho_U)$ is called a solution in the text above.
The uniqueness follows from Lemma \ref{lemma-uniqueness}.
If $X$ has a finite covering by elements of $\mathcal{B}$
(for example if $X$ is quasi-compact), then the theorem
is a consequence of Lemma \ref{lemma-solution-in-finite-case}.
In the general case we argue in exactly the same manner,
using transfinite induction and
Lemma \ref{lemma-glueing-increasing-union}.

\medskip\noindent
First we use transfinite recursion to choose opens $W_\alpha \subset X$
for any ordinal $\alpha$. Namely, we set $W_0 = \emptyset$.
If $\alpha = \beta + 1$ is a successor, then either $W_\beta = X$
and we set $W_\alpha = X$ or $W_\beta \not = X$ and we set
$W_\alpha = W_\beta \cup U_\alpha$ where
$U_\alpha \in \mathcal{B}$ is not contained in $W_\beta$.
If $\alpha$ is a limit ordinal we set
$W_\alpha = \bigcup_{\beta < \alpha} W_\beta$.
Then for large enough $\alpha$ we have $W_\alpha = X$.
Observe that for every $\alpha$ the open $W_\alpha$ is
a union of elements of $\mathcal{B}$. Hence if
$\mathcal{B}_\alpha = \{U \in \mathcal{B}, U \subset W_\alpha\}$, then
$$
S_\alpha = (\{K_U\}_{U \in \mathcal{B}_\alpha},
\{\rho_V^U\}_{V \subset U\text{ with }U, V \in \mathcal{B}_\alpha})
$$
is a system as in Lemma \ref{lemma-uniqueness} on the ringed space $W_\alpha$.

\medskip\noindent
We will show by transfinite induction that for every $\alpha$
the system $S_\alpha$ has a solution. This will prove the theorem
as this system is the system given in the theorem for large $\alpha$.

\medskip\noindent
The case where $\alpha = \beta + 1$ is a successor ordinal.
(This case was already treated in the proof of the lemma above
but for clarity we repeat the argument.)
Recall that $W_\alpha = W_\beta \cup U_\alpha$ for some
$U_\alpha \in \mathcal{B}$ in this case.
By induction hypothesis we have a solution
$(K_{W_\beta}, \{\rho^{W_\beta}_U\}_{U \in \mathcal{B}_\beta})$
for the system $S_\beta$.
Then we can consider the collection
$\mathcal{B}_\alpha^* = \mathcal{B}_\alpha \cup \{W_\beta\}$
of opens of $W_\alpha$ and we see that we obtain a system
$(\{K_U\}_{U \in \mathcal{B}_\alpha^*},
\{\rho_V^U\}_{V \subset U\text{ with }U, V \in \mathcal{B}_\alpha^*})$.
Note that this new system also satisfies condition (3)
by Lemma \ref{lemma-uniqueness} applied to the solution $K_{W_\beta}$.
For this system we have $W_\alpha = W_\beta \cup U_\alpha$.
This reduces us to the case handled in
Lemma \ref{lemma-solution-in-finite-case}.

\medskip\noindent
The case where $\alpha$ is a limit ordinal. Recall that
$W_\alpha = \bigcup_{\beta < \alpha} W_\beta$ in this case.
For $\beta < \alpha$ let
$(K_{W_\beta}, \{\rho^{W_\beta}_U\}_{U \in \mathcal{B}_\beta})$
be the solution for $S_\beta$.
For $\gamma < \beta < \alpha$ the restriction
$K_{W_\beta}|_{W_\gamma}$ endowed with the maps
$\rho^{W_\beta}_U$, $U \in \mathcal{B}_\gamma$
is a solution for $S_\gamma$. By uniqueness we get unique isomorphisms
$\rho_{W_\gamma}^{W_\beta} : K_{W_\beta}|_{W_\gamma} \to K_{W_\gamma}$
compatible with the maps $\rho^{W_\beta}_U$ and $\rho^{W_\gamma}_U$
for $U \in \mathcal{B}_\gamma$. These maps compose in the correct manner,
i.e., $\rho_{W_\delta}^{W_\gamma} \circ \rho_{W_\gamma}^{W_\beta}|_{W_\delta}
= \rho^{W_\delta}_{W_\beta}$ for $\delta < \gamma < \beta < \alpha$.
Thus we may apply Lemma \ref{lemma-glueing-increasing-union}
(note that the vanishing of negative exts is true for
$K_{W_\beta}$ by Lemma \ref{lemma-uniqueness} applied
to the solution $K_{W_\beta}$)
to obtain $K_{W_\alpha}$ and isomorphisms
$$
\rho_{W_\beta}^{W_\alpha} :
K_{W_\alpha}|_{W_\beta}
\longrightarrow
K_{W_\beta}
$$
compatible with the maps $\rho_{W_\gamma}^{W_\beta}$ for
$\gamma < \beta < \alpha$.

\medskip\noindent
To show that $K_{W_\alpha}$ is a solution we still need to construct the
isomorphisms $\rho_U^{W_\alpha} : K_{W_\alpha}|_U \to K_U$ for
$U \in \mathcal{B}_\alpha$ satisfying certain compatibilities.
We choose $\rho_U^{W_\alpha}$ to be the unique map such that
for any $\beta < \alpha$ and any $V \in \mathcal{B}_\beta$
with $V \subset U$ the diagram
$$
\xymatrix{
K_{W_\alpha}|_V \ar[r]_{\rho_U^{W_\alpha}|_V}
\ar[d]_{\rho_{W_\beta}^{W_\alpha}|_V}
& K_U|_V \ar[d]^{\rho_U^V} \\
K_{W_\beta} \ar[r]^{\rho_V^{W_\beta}}
& K_V
}
$$
commutes. This makes sense because
$$
(\{K_V\}_{V \subset U, V \in \mathcal{B}_\beta\text{ for some }\beta < \alpha},
\{\rho_V^{V'}\}_{V \subset V'\text{ with }V, V' \subset U
\text{ and }V, V' \in \mathcal{B}_\beta\text{ for some }\beta < \alpha})
$$
is a system as in Lemma \ref{lemma-uniqueness} on the ringed space $U$
and because $(K_U, \rho^U_V)$ and
$(K_{W_\alpha}|_U,  \rho_V^{W_\beta}\circ \rho_{W_\beta}^{W_\alpha}|_V)$
are both solutions for this system. This gives existence and uniqueness.
We omit the proof that these
maps satisfy the desired compatibilities (it is just bookkeeping).
\end{proof}

\end{document}
